\newcommand{\R}{\mathbb{R}}
\newcommand{\E}{\mathbb{E}}
\newcommand{\hvec}{\mathbf{h}}
\newcommand{\wvec}{\mathbf{w}}
\newcommand{\vvec}{\mathbf{v}}
\newcommand{\DeltaEff}{\Delta_{\mathrm{Eff}}}
\DeclareMathOperator*{\argmaxop}{arg\,max}
% One type size for every table. \fitwidth only ever shrinks, so a table wider than the line was
% set small and a narrow one was left at body size: across the paper the rendered table text ran from
% 7.1 pt to 10.0 pt, and the three ownership tables sat at the top of that range. At 8 pt with 3 pt
% column separation the widest table measures 453 pt against a 455 pt line, so no table is scaled at
% all and they all print at the same size. 8 pt is the largest size for which that holds.
\providecommand{\cftabsize}{\fontsize{8}{9.6}\selectfont}

% Shrink a table to the line, and to the page when it is the height that does not fit --
% table_cf_attr is 58 rows and overflowed a float page by 3.6pt once its width was fitted.
% max width/max totalheight only ever shrink, so a narrow table is left at body size.
\providecommand{\fitwidth}[1]{\adjustbox{max width=\linewidth,max totalheight=0.93\textheight}{#1}}
